\subsection{Measures, Charges, and Basic Examples}

\begin{definition} \label{def:measure-measurability-core-source-todo}\label{def:ring-sigma-ring-set-function} \hlt{Rings, $\sigma$-Rings, and Set Functions}\\
Let $\mathcal{A}\subseteq\mathcal{P}(X)$.
\begin{enumerate}[label=(\roman*)]
\setlength{\itemsep}{0pt}
\item $\mathcal{A}$ is a \hlt{ring of sets} iff $\emptyset\in\mathcal{A}$ and $A,B\in\mathcal{A}\Rightarrow A\cup B,\ A\setminus B\in\mathcal{A}$
\item $\mathcal{A}$ is a \hlt{$\sigma$-ring} iff it is a ring and $(A_n)\subseteq\mathcal{A}\Rightarrow\bigcup_{n=1}^{\infty}A_n\in\mathcal{A}$
\item A \hlt{set function} on $\mathcal{A}$ is $\phi:\mathcal{A}\to[-\infty,+\infty]$ with $\{+\infty,-\infty\}\nsubseteq\phi(\mathcal{A})$.
\item $\phi$ is \hlt{additive} iff $A\cap B=\emptyset\Rightarrow\phi(A\cup B)=\phi(A)+\phi(B)$
\item $\phi$ is \hlt{countably additive} iff $(A_n)\subseteq\mathcal{A}$ pairwise disjoint, $\bigcup_nA_n\in\mathcal{A}$ $\Rightarrow$ $\phi\left(\bigcup_{n=1}^{\infty}A_n\right)=\sum_{n=1}^{\infty}\phi(A_n)$.
\end{enumerate}
\hlt{Structural Commentary}: Rings support finite operations; $\sigma$-rings support countable unions.
\end{definition}

\begin{definition} \label{def:abstract-measure-on-measurable-space} \hlt{Measures on Measurable Spaces}\\
Let $(X,\mathcal{M})$ be a measurable space. A \hlt{measure} on $(X,\mathcal{M})$ is a map $\mu:\mathcal{M}\to[0,+\infty]$ such that:
\begin{enumerate}[label=(\roman*)]
\setlength{\itemsep}{0pt}
\item $\mu(\emptyset)=0$;
\item if $(E_n)\subseteq\mathcal{M}$ is pairwise disjoint, then $\mu\left(\bigcup_{n=1}^{\infty}E_n\right)=\sum_{n=1}^{\infty}\mu(E_n)$
\end{enumerate}
The triple $(X,\mathcal{M},\mu)$ is a \hlt{measure space}.\\
\hlt{Key Insight}: Measure is countably additive size on the chosen observable sets.
\end{definition}

\begin{definition} \label{def:abstract-measures-finite-sigma-finite} \hlt{Finite and $\sigma$-Finite Measures}\\
Let $(X,\mathcal{M},\mu)$ be a measure space.
\begin{enumerate}[label=(\roman*)]
\setlength{\itemsep}{0pt}
\item For nonnegative terms, $\sum_{n=1}^{\infty}a_n\defeq\lim_{N\to\infty}\sum_{n=1}^{N}a_n\in[0,+\infty]$.
\item $\mu$ is \hlt{finite} iff $\mu(E)<+\infty$ for every $E\in\mathcal{M}$.
\item $\mu$ is \hlt{$\sigma$-finite} iff there exist $E_n\in\mathcal{M}$ such that $X=\bigcup_{n=1}^{\infty}E_n$, $\mu(E_n)<+\infty$ for $n\in\N$
\end{enumerate}
\end{definition}

\begin{definition} \label{def:charges-finite-signed-measures} \hlt{Charges}\\
Let $(X,\mathcal{M})$ be a measurable space. A \hlt{charge} is a map $\lambda:\mathcal{M}\to\R$ with $\lambda(\emptyset)=0$ such that for every pairwise disjoint sequence $(E_n)\subseteq\mathcal{M}$, the series $\sum_{n=1}^{\infty}\lambda(E_n)$ converges in $\R$ and $\lambda(\bigcup_{n=1}^{\infty}E_n)=\sum_{n=1}^{\infty}\lambda(E_n)$.\\
\hlt{Key Insight}: Charges are real-valued signed measures with unconditional sums.
\end{definition}

\begin{example} \label{ex:first-abstract-measures} \hlt{First Abstract Measures}
\begin{enumerate}[label=(\roman*)]
\setlength{\itemsep}{0pt}
\item The zero set function $\mu(E)=0$ on $\mathcal{M}$ is a finite measure.
\item If $X\neq\emptyset$ and $\mathcal{M}=2^X$, $\mu(\emptyset)=0$ and $\mu(E)=+\infty$ for $E\neq\emptyset$ give a measure, neither finite nor $\sigma$-finite.
\item For $p\in X$, the \hlt{unit mass} $\delta_p(E)\defeq\mathbf{1}_E(p)$ on $2^X$ is a finite measure.
\item On $2^\N$, \hlt{counting measure} is $\#(E)$ for finite $E$ and $+\infty$ otherwise; it is $\sigma$-finite.\\
More generally, $(a_n)\subseteq[0,+\infty]$ gives $\mu(E)\defeq\sum_{n\in E}a_n$, and every measure on $2^\N$ is uniquely of this form.
\item If $X$ is uncountable and $\mathcal{M}=2^X$, $\mu(E)=0$ for countable $E$ and $\mu(E)=+\infty$ otherwise defines a measure.
\item On $2^\N$, $\mu(E)=0$ for finite $E$ and $\mu(E)=+\infty$ otherwise is not a measure, since $\N=\bigcup_{n=1}^{\infty}\{n\}$ gives $+\infty\neq0$; the constant $+\infty$ rule fails at $\emptyset$.
\item Lebesgue measure on $\mathcal{B}(\R)$ is $\sigma$-finite and gives length $b-a$ to standard intervals; Borel--Stieltjes measures come from increasing right-continuous functions as in Theorem \ref{thm:borel-stieltjes-measure-generation}.
\end{enumerate}
\hlt{Structural Commentary}: Measures may be zero, atomic, counting-type, cardinality-based, or geometric.\\
\hlt{Key Insight}: Countable additivity turns disjoint decompositions into nonnegative series.
\end{example}

\begin{proposition} \label{prop:elementary-measure-constructions} \hlt{Elementary Measure Constructions}\\
Let $(X,\mathcal{M})$ be a measurable space.
\begin{enumerate}[label=(\roman*)]
\setlength{\itemsep}{0pt}
\item If $\mu$ is a measure on $\mathcal{M}$ and $A\in\mathcal{M}$, then $\mu_A(E)\defeq\mu(A\cap E)$ defines a measure $\mu_A$ on $\mathcal{M}$.
\item If $\mu_1,\ldots,\mu_m$ are measures on $\mathcal{M}$ and $a_1,\ldots,a_m\geq0$, then $\lambda\defeq\sum_{j=1}^{m}a_j\mu_j$ is a measure on $\mathcal{M}$.
\item If $(\mu_n)_{n\geq1}$ are measures on $\mathcal{M}$ with $\mu_n(X)=1$, then $\lambda\defeq\sum_{n=1}^{\infty}2^{-n}\mu_n$ is a measure on $\mathcal{M}$ and $\lambda(X)=1$.
\end{enumerate}
\begin{proof}
\begin{enumerate}[label=(\roman*)]
\setlength{\itemsep}{0pt}
\item $\mu_A(\emptyset)=0$. If $(E_k)$ is disjoint in $\mathcal{M}$, then $A\cap\bigsqcup_kE_k=\bigsqcup_k(A\cap E_k)$, whence $\mu_A(\bigsqcup_kE_k)=\sum_k\mu_A(E_k)$.
\item $\lambda(\emptyset)=0$. For disjoint $(E_k)$, finite summation gives $\lambda(\bigsqcup_kE_k)=\sum_{j=1}^{m}a_j\sum_k\mu_j(E_k)=\sum_k\lambda(E_k)$.
\item $\lambda(\emptyset)=0$. For disjoint $(E_k)$, nonnegative reordering gives $\lambda(\bigsqcup_kE_k)=\sum_n2^{-n}\sum_k\mu_n(E_k)=\sum_k\lambda(E_k)$.\\
Also $\lambda(X)=\sum_n2^{-n}\mu_n(X)=\sum_n2^{-n}=1$.
\end{enumerate}
\hlt{Structural Commentary}: New measures come from restriction, superposition, and countable convex mixtures.\\
\hlt{Key Insight}: Nonnegative series make the bookkeeping order-independent.
\end{proof}
\end{proposition}

\begin{proposition} \label{prop:additive-set-function-basic-laws} \hlt{Basic Laws for Additive Set Functions}\\
Let $\phi$ be additive on a ring $\mathcal{A}$. Then:
\begin{enumerate}[label=(\roman*)]
\setlength{\itemsep}{0pt}
\item $\phi(\emptyset)=0$;
\item $\phi(A_1\cup\cdots\cup A_n)=\phi(A_1)+\cdots+\phi(A_n)$ for pairwise disjoint $A_i\in\mathcal{A}$;
\item $\phi(A\cup B)+\phi(A\cap B)=\phi(A)+\phi(B)$;
\item if $\phi$ is nonnegative and $A\subseteq B$, then $\phi(A)\leq\phi(B)$;
\item if $B\subseteq A$ and $\abs{\phi(B)}<+\infty$, then $\phi(A\setminus B)=\phi(A)-\phi(B)$.
\end{enumerate}
\begin{proof}
\begin{enumerate}[label=(\roman*)]
\setlength{\itemsep}{0pt}
\item Apply additivity to the empty disjoint union.
\item Induct on $n$, applying additivity to $(A_1\cup\cdots\cup A_{n-1})\cup A_n$.
\item By disjoint decompositions, $A\cup B=(A\setminus B)\cup(B\setminus A)\cup(A\cap B)$; $A=(A\setminus B)\cup(A\cap B)$; $B=(B\setminus A)\cup(A\cap B)$
\item If $A\subseteq B$, then $B=A\cup(B\setminus A)$ disjointly, so nonnegativity gives $\phi(A)\leq\phi(B)$.
\item If $B\subseteq A$, then $A=B\cup(A\setminus B)$ disjointly; since $\abs{\phi(B)}<+\infty$, subtract $\phi(B)$.
\end{enumerate}
\hlt{Structural Commentary}: Finite additivity plus disjoint decompositions controls all elementary identities.\\
\hlt{Key Insight}: Rewrite overlaps as disjoint pieces before applying $\phi$.
\end{proof}
\end{proposition}

\begin{theorem} \label{thm:continuity-from-below-set-functions} \hlt{Continuity from Below}\\
Suppose $\phi$ is countably additive on a ring $\mathcal{A}$.\\
If $A_1\subseteq A_2\subseteq A_3\subseteq\cdots$, and $A=\bigcup\limits_{n=1}^{\infty}A_n$ with all $A_n,A\in\mathcal{A}$, then $\phi(A_n)\to\phi(A)$.
\begin{proof}
Put $B_1=A_1$ and $B_n=A_n\setminus A_{n-1}$ for $n\geq2$. Then $B_n$ are pairwise disjoint, $A_n=B_1\cup\cdots\cup B_n$, and $A=\bigcup_nB_n$. Countable additivity turns $\phi(A_n)$ and $\phi(A)$ into partial sums and full sum of same series.\\
\hlt{Structural Commentary}: Increasing unions are disjoint tail decompositions in disguise.\\
\hlt{Key Insight}: Countable additivity converts set convergence into series convergence.
\end{proof}
\end{theorem}

\begin{definition} \label{def:charges-and-variation-measures} \hlt{Variation Candidates of a Charge}\\
Let $\lambda$ be a charge on $\mathcal{M}$. For $E\in\mathcal{M}$, define
\begin{enumerate}[label=(\roman*)]
\setlength{\itemsep}{0pt}
\item $\mathfrak{P}_{\mathcal{M}}(E)\defeq\{(A_j)_{j=1}^{n}:n\in\N,\ E=\bigsqcup_{j=1}^{n}A_j,\ A_j\in\mathcal{M}\}$, the finite measurable partitions of $E$;
\item $\pi_\lambda(E)\defeq\sup\{\lambda(A):A\subseteq E,\ A\in\mathcal{M}\}$, the \hlt{positive variation} candidate;
\item $\nu_\lambda(E)\defeq\sup\{\sum_{j=1}^{n}\abs{\lambda(A_j)}:(A_j)_{j=1}^{n}\in\mathfrak{P}_{\mathcal{M}}(E)\}$, the \hlt{total variation} candidate.
\end{enumerate}
For $E=\emptyset$, the empty partition is allowed and contributes $0$.
\end{definition}

\begin{proposition} \label{prop:charges-linear-combinations-and-variation} \hlt{Charge Operations and Variation Measures}\\
Let $\lambda,\lambda_1,\ldots,\lambda_m$ be charges on $\mathcal{M}$.
\begin{enumerate}[label=(\roman*)]
\setlength{\itemsep}{0pt}
\item If $a_1,\ldots,a_m\in\R$, then $\Lambda\defeq\sum_{j=1}^{m}a_j\lambda_j$ is a charge.
\item If $f\in\mathscr{L}^1(\mu)$ on $(X,\mathcal{M},\mu)$, then $\lambda_f(A)\defeq\int_A f\,d\mu$ is a charge.
\item The set functions $\pi_\lambda$ and $\nu_\lambda$ are measures on $\mathcal{M}$.
\end{enumerate}
\begin{proof}
\begin{enumerate}[label=(\roman*)]
\setlength{\itemsep}{0pt}
\item If $E=\bigsqcup_nE_n$, then $\Lambda(E)=\sum_{j=1}^{m}a_j\sum_n\lambda_j(E_n)=\sum_n\Lambda(E_n)$.
\item By Theorem \ref{thm:integral-countably-additive-set-function}, $A\mapsto\int_A f\,d\mu$ is countably additive, and $\abs{\lambda_f(A)}\leq\int_X\abs{f}\,d\mu<+\infty$.
\item Let $E=\bigsqcup_nE_n$. For $A\subseteq E$, $A=\bigsqcup_n(A\cap E_n)$, so $\lambda(A)\leq\sum_n\pi_\lambda(E_n)$; by suprema, $\pi_\lambda(E)\leq\sum_n\pi_\lambda(E_n)$. Conversely, for $N\in\N$ choose $A_n\subseteq E_n$ with $\lambda(A_n)>\pi_\lambda(E_n)-\varepsilon2^{-n}$ for $n\leq N$.\\
Then $\pi_\lambda(E)\geq\sum_{n=1}^{N}\pi_\lambda(E_n)-\varepsilon$. Let $N\to\infty$ and $\varepsilon\downarrow0$. For $\nu_\lambda$, every $(A_j)_{j=1}^{r}\in\mathfrak{P}_{\mathcal{M}}(E)$ satisfies $\sum_{j=1}^{r}\abs{\lambda(A_j)}\leq\sum_n\nu_\lambda(E_n)$ after refining by $A_j\cap E_n$, hence $\nu_\lambda(E)\leq\sum_n\nu_\lambda(E_n)$. Conversely, combine near-optimal finite partitions of $E_1,\ldots,E_N$, add $\bigcup_{n>N}E_n$ as one part; let $N\to\infty$.
\end{enumerate}
\hlt{Structural Commentary}: Variation turns signed additivity into measure.\\
\hlt{Key Insight}: Refine partitions against a disjoint decomposition.
\end{proof}
\end{proposition}
